% ============ 前言部分 (原书 pp.4-6) ============
% 记号与约定 / 预备知识与参考文献 / 致谢 / 评注
% (原书目录不收录这些条目, 故用 \chapter* 不加入目录)

\chapter*{记号与约定}

关于簇的约定与我的\emph{代数几何}讲义中的相同.例如，域 $k$ 上的仿射代数(affine algebra)是有限生成的 $k$-代数 $A$，使得对 $k$ 的一个(从而每个)代数闭包 $k^{\mathrm{al}}$，$A \otimes_k k^{\mathrm{al}}$ 无非零幂零元---这蕴含 $A$ 本身无幂零元.对这样的 $k$-代数，我们联系一个环化空间 $\operatorname{Specm}(A)$(赋予 $k$-代数层结构的拓扑空间)，而 $k$ 上的仿射簇是同构于这种形式的环化空间.$k$ 上的簇是环化空间 $(X, \mathcal{O}_X)$，它有一个有限开覆盖 $X = \bigcup U_i$，使得每个 $(U_i, \mathcal{O}_X|_{U_i})$ 都是仿射簇，并且满足分离公理.我们常用 $X$ 同时表示 $(X, \mathcal{O}_X)$ 及其底层拓扑空间.簇的正则映射有时也称为簇的态射.

对偏好概形的读者：簇是域 $k$ 上有限型的分离，几何既约概形 $X$，其中非闭点被去掉.

若 $X$ 是 $k$ 上的簇，$K \supset k$，则 $X(K)$ 是 $X$ 中坐标在 $K$ 里的点的集合，而 $X_K$ 或 $X_{/K}$ 是由 $X$ 得到的 $K$ 上簇\footnote{我们的术语遵循 Grothendieck.存在一种相冲突的术语，基于 Weil 的 \emph{Foundations}，常被代数群理论的研究者使用，它把两者反过来写.}.例如，若 $X = \operatorname{Specm}(A)$，则 $X(K) = \operatorname{Hom}_{k\text{-alg}}(A, K)$，且 $X_K = \operatorname{Specm}(A \otimes_k K)$.一般地，$X(K)$ 只是一个集合，但我通常赋予 $X(\mathbb{C})$ 自然的复拓扑.域 $k$ 的可分闭包 $k^{\mathrm{sep}}$ 是 $k$ 的代数扩张，使得每个系数在 $k$ 中的可分多项式在 $k^{\mathrm{sep}}$ 中都有根.

关于概形的术语是标准的，只有一点除外：我总假定我们的环是 Noether 的，我们的概形是局部 Noether 的.

关于环的术语是标准的.特别地，环的同态 $A \to B$ 把 1 映到 1.环同态 $A \to B$ 是有限的，而 $B$ 是有限 $A$-代数，如果 $B$ 作为 $A$-模有限生成.当 $A$ 是局部环时，我常用 $\mathfrak{m}_A$ 记其(唯一的)极大理想.局部环的局部同态是同态 $f \colon A \to B$，使得 $f^{-1}(\mathfrak{m}_B) = \mathfrak{m}_A$(等价地，$f(\mathfrak{m}_A) \subset \mathfrak{m}_B$).

一般地，当我从拓扑空间上的层论汲取动机时，我假定空间是 Hausdorff 的，即不是那种有点的闭包为全空间的怪空间.

\begin{itemize}
  \item $X \subset Y$：$X$ 是 $Y$ 的子集(不一定是真子集);
  \item $X \overset{\mathrm{def}}{=} Y$：$X$ 定义为 $Y$，或按定义等于 $Y$；
  \item $X \cong Y$：$X$ 同构于 $Y$；
  \item $X \approx Y$：$X$ 与 $Y$ 典范同构(或存在给定的或唯一的同构).
\end{itemize}

\chapter*{预备知识与参考文献}

\textbf{同调代数} 我假定读者熟悉 Abel 范畴与导出函子的语言.这些主题在我的\emph{类域论}讲义 pp.69--76 中有概述，若干书中有完整叙述，例如
\begin{quote}
  Weibel, C.A., \emph{An Introduction to Homological Algebra}\footnote{在第 9 页，Weibel 定义 $C[p]_n = C_{n-p}$.正确的原始定义(代数几何学家与算术几何学家普遍使用)是 $C[p]_n = C_{n+p}$(见 Hartshorne, R., \emph{Residues and Duality}, 1966, p.26).此外，在 Weibel 中，``函子范畴''不一定是范畴.}, Cambridge U.P., 1994.
\end{quote}
我们不能避免使用谱序列---谱序列的简短概述见我的\emph{平展上同调}一书 pp.307--309，完整处理见 Weibel 的书第 5 章.

\textbf{层论} 平展上同调以通常拓扑意义下的层上同调理论为模型.本讲义中的许多材料与例如下面的材料平行：
\begin{quote}
  Iversen, B., \emph{Cohomology of Sheaves}, Springer, 1986.
\end{quote}

\textbf{代数几何} 我假定读者熟悉代数簇的理论，例如我的\emph{代数几何}讲义(Math. 631).此外，有时我会提到概形，因此读者应熟悉概形的基本语言，例如 Hartshorne, \emph{Algebraic Geometry}, Springer 1977 第 II 章的前 3 节，Eisenbud and Harris, \emph{Schemes}, Wadsworth, 1992 的第 1 章，或 Shafarevich, \emph{Basic Algebraic Geometry}, 2nd Edition, Springer, 1994 的第 V 章.

交换代数我通常参考
\begin{quote}
  Atiyah, M., and MacDonald, I., \emph{Introduction to Commutative Algebra}, Addison-Wesley, 1969.
\end{quote}

\textbf{平展上同调} 有下列书籍：
\begin{quote}
  Freitag, E., and Kiehl, R., \emph{Etale Cohomology and the Weil Conjecture}, Springer, 1988.\\
  Milne, J., \emph{Etale Cohomology}, Princeton U.P. 1980 (引用为 EC).\\
  Tamme, \emph{Introduction to Etale Cohomology}, Springer.
\end{quote}

原始文献：
\begin{quote}
  Artin, M., \emph{Grothendieck Topologies}, Lecture Notes, Harvard University Math. Dept. 1962.\\
  Artin, M., \emph{Th\'eor\`emes de Repr\'esentabilit\'e pour les Espaces Alg\'ebriques}, Presses de l'Universit\'e de Montr\'eal, Montr\'eal, 1973\\
  Grothendieck, A., et al., \emph{S\'eminaire de G\'eom\'etrie Alg\'ebrique}.\\
  SGA 4 (with Artin, M., and Verdier, J.-L.). \emph{Th\'eorie des topos et cohomologie \'etale des sch\'emas} (1963--64). Springer Lecture Notes 1972--73.\\
  SGA $4\frac{1}{2}$ (by Deligne, P., with Boutot, J.-F., Illusie, L., and Verdier, J.-L.) \emph{Cohomologie \'etale}. Springer Lecture Notes 1977.\\
  SGA 5 \emph{Cohomologie l-adique et fonctions L} (1965--66). Springer Lecture Notes 1977.\\
  SGA 7 (with Deligne, P., and Katz, N.) \emph{Groupes de monodromie en g\'eom\'etrie alg\'ebriques} (1967--68). Springer Lecture Notes 1972--73.
\end{quote}

除 SGA $4\frac{1}{2}$ 外，这些都是 Grothendieck 在 I.H.E.S. 领导的著名讨论班.

我引用我的如下讲义.
\begin{quote}
  FT \emph{Field and Galois Theory} (v4.40 2013).\\
  AG \emph{Algebraic Geometry} (v5.22 2012).\\
  AV \emph{Abelian Varieties} (v2.00 2008).\\
  CFT \emph{Class Field Theory} (v4.01 2011).
\end{quote}

\chapter*{致谢}

我感谢以下人士为这些讲义的早期版本提供勘误与意见：Martin Bright; Rex Cheung; Sungmun Cho; Ming-chang Kang; Michiel Kosters; Akhil Matthew; Carl Mautner and others at UT Austin; Behrooz Mirzaii; Eric Moorhouse; Sean Rostami; Steven Spallone; Immanuel Stampfli; Sun Shenghao; Zach Teitler; Bhupendra Nath Tiwar; Ravi Vakil.

\chapter*{评注}

平展上同调中的主要定理都在 SGA 4 和 SGA 5 中证明，但常常是在不必要地严格的假设之下.其中一些假设后来被去除了，但所用的证明大量使用了那些讨论班的内容.因此这门学科的结构需要重新思考.此外，代数空间应当更充分地融入这门学科(见 Artin 1973).de Jong 的分解定理(\emph{Smoothness, semi-stability and alterations}. Inst. Hautes \'Etudes Sci. Publ. Math. No. 83 (1996), 51--93)很可能允许许多改进.最后，诸如相交上同调与 Borel-Moore 同调这样的主题需要补充到讲述中.这些都不会在这份讲义中尝试.

\begin{quote}
  我认为，Grothendieck 在 1958 年 Serre 于 Chevalley 讨论班上的报告之后之所以确信平展局部化会给出正确的 $H^i$，是因为一旦你有了曲线的正确上同调，那么通过曲线纤维化与 d\'evissage，你也应当能达到更高的 $H^i$.(Illusie, NAMS 2010.)
\end{quote}

\begin{quote}
  Grothendieck 第一次来哈佛是在 1958 年.Artin：是的.那时我还不知道概形是什么.但第二次，在 '61 年，我听说他有平展上同调的想法，当他到达时我问他我是否可以思考它，那就是开始.他当时并没有在做这个---他有这个想法但把它搁置了.他直到我证明了第一个定理才开始做它.他极其活跃，但在那些年里，这可能是他唯一一件没有立刻做的事，原因不明.那年秋天我思考了它，我们争论定义应该是什么[笑].然后我在 '62 年秋天做了一个讨论班.(\emph{Recountings: conversations with MIT mathematicians}, p.358.)
\end{quote}
